\chapter{Studio della Complessità}

In questo capitolo viene analizzata la complessità temporale e spaziale 
dell'algoritmo implementato per il problema dei Percorsi nell'Arcipelago.

\section{Complessità Temporale}

\subsection{Lettura dell'Input}

La lettura del file di input richiede la lettura di $P$ tratte, ciascuna 
in tempo costante. La complessità è:

\begin{equation}
    T_{input} = O(P)
\end{equation}

\subsection{Costruzione del Grafo}

Il metodo \texttt{addEdge} viene invocato $P$ volte, ciascuna in tempo 
$O(1)$ (inserimento in fondo a un vettore). La complessità è:

\begin{equation}
    T_{build} = O(P)
\end{equation}

\subsection{Variante di Dijkstra}

L'algoritmo utilizza una coda di priorità implementata come min-heap 
tramite \texttt{std::priority\_queue}. Ogni nodo può essere inserito 
nella coda più volte (una per ogni arco che lo raggiunge), quindi il 
numero massimo di elementi nella coda è $O(P)$.

Le operazioni principali sono:

\begin{itemize}
    \item \textbf{Estrazione dalla coda}: al più $O(P)$ estrazioni, 
    ciascuna in tempo $O(\log P)$;
    \item \textbf{Inserimento nella coda}: al più $O(P)$ inserimenti, 
    ciascuno in tempo $O(\log P)$;
    \item \textbf{Rilassamento degli archi}: ogni arco viene esaminato 
    al più due volte (grafo non orientato), per un totale di $O(P)$ 
    rilassamenti in tempo $O(1)$ ciascuno.
\end{itemize}

La complessità complessiva della variante di Dijkstra è:

\begin{equation}
    T_{dijkstra} = O(P \log P)
\end{equation}

\subsection{Ricostruzione dei Percorsi}

Per ogni città $v$ da $1$ a $N$, il metodo \texttt{reconstructPath} 
risale il vettore \texttt{parent} per al più $N$ passi. La complessità 
complessiva per tutti i percorsi è:

\begin{equation}
    T_{reconstruct} = O(N^2)
\end{equation}

Nel caso pratico, la lunghezza media dei percorsi è molto inferiore 
a $N$, quindi questa stima è pessimistica.

\subsection{Complessità Totale}

La complessità temporale totale dell'algoritmo è:

\begin{equation}
    T_{tot} = O(P \log P + N^2)
\end{equation}

Poiché $P \leq 10000$ e $N \leq 1000$, entrambi i termini sono 
dominati da valori molto contenuti nella pratica.

\section{Complessità Spaziale}

Le strutture dati utilizzate occupano la seguente memoria:

\begin{itemize}
    \item \textbf{Lista di adiacenza} \texttt{adj}: $O(N + P)$ per 
    memorizzare tutti i nodi e le tratte;
    \item \textbf{Vettori ausiliari} \texttt{dist}, \texttt{totalCost}, 
    \texttt{maxEdge}, \texttt{parent}: $O(N)$ ciascuno, per un totale 
    di $O(N)$;
    \item \textbf{Coda di priorità}: al più $O(P)$ elementi 
    simultaneamente;
    \item \textbf{Vettore \texttt{path}} in \texttt{reconstructPath}: 
    $O(N)$ per il percorso più lungo possibile.
\end{itemize}

La complessità spaziale totale è:

\begin{equation}
    S = O(N + P)
\end{equation}

\section{Confronto con Dijkstra Classico}

\begin{center}
\begin{tabular}{|l|c|c|}
\hline
\textbf{Algoritmo} & \textbf{Complessità Temporale} & \textbf{Complessità Spaziale} \\
\hline
Dijkstra classico      & $O(P \log P)$ & $O(N + P)$ \\
Variante con sconto    & $O(P \log P)$ & $O(N + P)$ \\
\hline
\end{tabular}
\end{center}

\vspace{0.5cm}
La variante proposta ha la stessa complessità asintotica del Dijkstra 
classico. L'unica differenza è che ogni stato nella coda contiene 
tre campi invece di due (\texttt{totalCost}, \texttt{maxCost}, 
\texttt{node}), con un overhead costante sia in tempo che in spazio.